\section{Jankurai Standard and Conformance}

Jankurai is the control plane. It is not a general-purpose claim that all AI repositories should use one stack or one tool. It is a versioned standard for repositories that want to claim agent-native conformance. The standard binds prose, machine-readable maps, generated-zone policy, audit output, CI behavior, and repair evidence.

The adoption loop is itself part of the standard: `jankurai init` installs the control plane, `jankurai doctor` reports drift, and `jankurai audit` produces the proof ledger. A repository that only mentions the standard but does not install or execute it is not conformant.

Adoption is progressive rather than binary. The install ladder is \texttt{agents} \(\rightarrow\) \texttt{score} \(\rightarrow\) \texttt{ci} \(\rightarrow\) \texttt{full} \(\rightarrow\) \texttt{ratchet}. \texttt{jankurai init --level agents} installs only the root router, the short standard, the master plan, and thin provider adapters. This is the smallest useful change because agent behavior is shaped before the repository claims anything about conformance. \texttt{--level score} adds local scoring manifests: owner map, test map, generated zones, proof lanes, audit policy, standard version, and minimal command recipes. \texttt{--level ci} adds observe-mode CI plus security policy stubs and evidence paths; it must not enforce score 85 by default. \texttt{--level full} preserves the profile-driven scaffold for teams that want docs, contracts, database placeholders, UX QA policy, security controls, and profile-specific structure in one step. Ratchet enforcement remains a separate CI operation after a baseline has been reviewed.

\begin{table}[t]
\caption{Progressive jankurai install levels.}
\label{tab:install-levels}
\centering
\small
\begin{tabularx}{\columnwidth}{L{0.18\columnwidth} Y}
\toprule
\textbf{Level} & \textbf{Installed surface} \\
\midrule
\texttt{agents} & Root instructions, standard brief, master plan, and provider adapters. No score claim. \\
\texttt{score} & Agents plus owner, proof, generated-zone, audit, version, and local recipe files for advisory scoring. \\
\texttt{ci} & Score plus observe-mode workflow, security policy, and stub evidence path. No default score gate. \\
\texttt{full} & Selected profile unchanged: docs, contracts, DB placeholders, UX, security, and CI where declared. \\
\texttt{ratchet} & Explicit CI install mode that compares against an accepted baseline and blocks regression. \\
\bottomrule
\end{tabularx}
\end{table}

\begin{table}[t]
\caption{Jankurai conformance levels.}
\label{tab:conformance}
\centering
\small
\begin{tabularx}{\columnwidth}{L{0.16\columnwidth} Y}
\toprule
\textbf{Level} & \textbf{Meaning} \\
\midrule
HL0 & Unscored: no reliable root instructions or proof lane. \\
HL1 & Advisory: audit runs and emits JSON/Markdown, but findings do not block merge. \\
HL2 & Guarded: critical caps, generated-zone mutation, and missing fast lane block merge. \\
HL3 & Standard: high and critical findings block merge; score floor and conformance metadata are enforced. \\
HL4 & Ratchet: score cannot regress without an exception; repair queue is tracked to closure. \\
HL5 & Release contract: audit, security, contracts, DB, e2e, and artifact versions gate release. \\
\bottomrule
\end{tabularx}
\end{table}

MUST, SHOULD, and MAY are used in the RFC sense inside the standard. A repository claiming HL3 or higher MUST provide a root agent router, owner map, test map, generated-zone manifest, version manifest, one-command fast lane, audit JSON, audit Markdown, and an ordered \texttt{agent\_fix\_queue}. It MUST NOT hand-edit generated zones, hide durable truth in the wrong layer, or merge high-risk changes without a mapped proof lane. It SHOULD provide SARIF output when the host CI can ingest it, plus JUnit, GitHub summaries, JSONL repair queues, and issue exports where they improve repair routing.

Stable rule identifiers make the standard repairable. The initial rule family covers dead markers, generated-zone mutation, owner and proof routing, Python truth, wrong-layer DB access, handwritten contracts, false-green risk, generated-code security, secret sprawl, hostile context, overbroad agency, rendered UX, accessibility, setup/context gaps, supply-chain drift, opaque observability, performance or concurrency drift, and streaming runtime drift. Appendix~A expands these IDs into required evidence and repair fields.

The audit loop is intentionally concrete. A failing repository might emit \texttt{HLT-010-SECRET-SPRAWL} for an exposed token in a fixture, \texttt{HLT-007-HANDWRITTEN-CONTRACT} for a handwritten web DTO, and \texttt{HLT-013-RENDERED-UX-GAP} for a checkout change without browser artifacts. The \texttt{agent\_fix\_queue} should order those by TLR priority, route them to \texttt{security}, \texttt{contract}, and \texttt{web} lanes, and require repair receipts: rotated credential and secret-scan output, regenerated client and drift check, screenshot/ARIA/geometry report for the changed viewport. Proof beats claims. A green score without those receipts is not conformance; it is merely an absence of local evidence.

Centerline drift is the score delta between a repository's claimed standard and its observed behavior. Hard caps are versioned policy, not final empirical truth. The point is to make governance explicit: teams can argue about weights and caps in releases, not in scattered PR comments.
